\section{Background and Related Work}
\label{sec:background}

\subsection{DABStep}
\label{sec:dabstep}

DABStep~\cite{dabstep} is a data-analysis benchmark built by Adyen and Hugging
Face around a payments warehouse: 138{,}236 transactions, a fee-rule table,
merchant and merchant-category-code (MCC) reference data, and two documents
of human-written context, \texttt{manual.md} (22{,}127 characters) and a
dataset README. Three of its properties shape this study.

\textbf{Its difficulty is semantic, not syntactic.} The hard split does not
require exotic SQL. It requires knowing the fee-matching rule: a \texttt{NULL}
in a \texttt{fees} column means \emph{this rule applies to every value of that
column}, not \emph{unknown}, and matching additionally involves list
membership over account type, authorisation characteristics indicator (ACI)
and merchant category code. An agent that reads the schema and writes
idiomatic SQL gets a clean, wrong number, which is the failure a context
layer claims to fix. Among text-to-SQL benchmarks this sets DABStep apart
from Spider~\cite{spider} and Spider~2.0~\cite{spider2}: its difficulty is
documentation, not schema scale.

\textbf{Golds are withheld and grading is centralised}, which raises the cost
of the overfitting that makes many text-to-SQL results hard to
trust~\cite{ladder,leaderboard-illusion}. It does not remove it: the sources
an agent reads are public, and Section~\ref{sec:golds} recovers 59
hard-split answers from them.

\textbf{Its question space is small.} DABStep expands 450 test tasks from a
smaller set of core questions by varying merchants, schemes, months and
thresholds, following GSM-Symbolic's parameterisation~\cite{gsm-symbolic}.
Normalising those varying entities away leaves 113 distinct question shapes,
of which 12 cover half the benchmark and 27 cover three quarters. In our task
set 294 of 332 hard tasks concern fees. Section~\ref{sec:threats} treats this
as the sharpest limit on the result, and Section~\ref{sec:prior} explains
why it favours pre-computed views.

The benchmark's own paper reports its strongest agent at 14.55\% on the
hard split against 76.39\% on the easy one~\cite{dabstep}. Hard-split
accuracies from credible named leaderboard entries since cluster around
20--26\% (Google's \texttt{simple\_baseline} at 26\%, an Adyen GPT-5.4 entry
at 22.2\%, a Hugging Face Claude~4 Sonnet entry at 19.8\%), and we use that
band as the validity check on our reimplementation of the prompt-based
approach (Section~\ref{sec:headline}).

\subsection{Three kinds of layer, and when their semantics are compiled}
\label{sec:layers}

Three families of approach exist and are routinely conflated. \emph{Semantic
layers} in the classical sense --- dbt metrics, Cube, LookML --- define
metrics as executable artifacts, so the agent calls a metric rather than
deriving it and correctness is inherited from the definition. \emph{Context
layers} store natural-language documentation next to the data and retrieve
it at query time; MotherDuck Guides~\cite{motherduck-guides} is the
best-documented instance. The agent still writes the SQL, but writes it
knowing what the columns mean. \emph{Data contracts} originate in governance
rather than assistance: a declarative artifact stating what a dataset
contains, who owns it and what may be done with it, with the Open Data
Contract Standard~\cite{odcs} the emerging specification. They inherit a
longer line of dataset documentation, from datasheets~\cite{datasheets} to
Croissant's metadata a tool loads~\cite{croissant}. MotherDuck's own primer
draws the first line the same way: a semantic layer holds definitions
``written as logic that compiles to SQL'' and a context layer ``holds
everything else''~\cite{motherduck-primer}. The database community's framing
of the wider problem is TAG~\cite{tag}.

The artifact under test is of the third kind and serves both purposes at
once: the same YAML supplies the semantics the agent reads and the rules the
tool layer enforces (Section~\ref{sec:system}). The distinction that matters
for this paper is \emph{when} the semantics are compiled
(Table~\ref{tab:compile}). A macro is compiled ahead of time by a human and
covers the metrics somebody anticipated; a contract is compiled per query by
the model and covers whatever question arrives.

\begin{table}[t]
\centering
\small
\begin{tabular}{@{}L{0.27\columnwidth}L{0.26\columnwidth}L{0.33\columnwidth}@{}}
\toprule
 & \tabhead{Views / macros} & \tabhead{Declarative contract} \\
\midrule
Compiled & ahead of time, by a human & per query, by the model \\
Covers & anticipated metrics & any question in the domain \\
Cost & one macro per metric, maintained & one description per domain \\
Fails when & nobody anticipated the question & the model misapplies the rule \\
\bottomrule
\end{tabular}
\caption{Two ways to deliver the same semantics. The compiled contract of
Section~\ref{sec:ceiling} is the left column built from the right column's
content, which is what lets this paper measure the difference between them.}
\label{tab:compile}
\end{table}

\subsection{What has been measured, and what this paper adds}
\label{sec:prior}

\textbf{Context layers on DABStep.} MotherDuck's Guides --- markdown context
stored in the warehouse and fetched through an MCP server --- raise DABStep
accuracy by 72 points and cut cost per run by 55\%, with Gemini~3 Flash
answering 418 of 419 held-out questions (99.8\%)~\cite{motherduck-guides}; a
companion report reaches 98.6\% with a locally hosted 27B model at 4-bit
quantization~\cite{motherduck-local}. A second MotherDuck report makes the
contribution of pre-computation legible because it reports a progression
rather than a number~\cite{motherduck-semantic}: retrieving context fragments
by vector search ``capped out around 88\%''; baking the knowledge into the
warehouse as schema comments, macros and derived tables reached 93\%; a
hierarchical semantic layer over the raw data, authored by a large model and
refined iteratively, reached 100\%. Each step moves work out of the agent and
into an artifact prepared in advance. The system atop the validated leaderboard, at 89.95\% hard, is built the
same way, by validating candidate helper functions against the benchmark's
gold answers~\cite{nvidia-kgmon}; nothing in the benchmark forbids that,
and we allege no violation. The second-placed system, at 87.57\%, is
described in a post we could not retrieve in full~\cite{oceanbase-datapilot}.

A reader will otherwise stop at the difference between those numbers and our
best hard-split figure of 77.4\%, so we say now what it is made of. Pre-built
views are maximally effective when the question space is enumerable in
advance, which is this benchmark's construction (Section~\ref{sec:dabstep}).
A macro contributes nothing to a question nobody wrote it for; the
declarative layer measured here is the layer that must cover the
unanticipated question. \textbf{Our hard-split figures measure a harder task,
not a worse attempt at the same one}, and Section~\ref{sec:ceiling} puts a
number on how much harder: the compiled form of our own contract answers
53\% of the hard split outright and none of the rest. The two layers are
complementary. Macros optimise the head of the question distribution, a
contract covers the tail, and in a system with both, the contract is what
tells the agent which macro exists and when to call it. We do not claim the
contract approach scales better, only that it degrades differently: a
missing macro degrades to nothing, a contract to a model reasoning from a
description.

The tuned artifacts also do not travel: the author of the 100\% result
writes that ``the tuned artifact is coupled'' to the serving
model~\cite{motherduck-semantic}, and independent commentary on Guides urges
measuring ``whether the Guide generalizes to new questions, not whether it
clears the benchmark you tuned it against''~\cite{corrdyn-guides}. Our
contract was authored from the vendor manual alone, frozen before any
question was read, and run unchanged against four model families, which is
why it can be published with the paper (\appref{app:leaderboard}).

\textbf{Methods this paper borrows.} Stripping the semantic content of a
prompt component while holding its form fixed, to see what the form alone
was worth, is the move of~\cite{min-demonstrations}; here the component is a
whole retrieval layer. Two results frame the delivery mechanism of
Section~\ref{sec:why}: retrieval quality falls for material buried
mid-prompt~\cite{lost-in-the-middle}, and what a tool's documentation says
changes what an agent does with it~\cite{tool-documentation}. For governance,
capability-based enforcement outside the model is the design behind
CaMeL~\cite{camel}, policy compliance under a prompt-stated document is what
$\tau$-bench measures~\cite{tau-bench}, and RuLES finds models breaking
simple stated rules with no adversary present~\cite{rules}, the published
expectation Section~\ref{sec:governance} runs against. Reporting cost beside
accuracy follows~\cite{agents-that-matter}. Paired testing across several
models is standard~\cite{semantic-paired,dbt-benchmark}, and the latter's 20
repetitions per question is better variance discipline than ours
(Section~\ref{sec:threats}); McNemar's test for learners executed once is
Dietterich's recommendation~\cite{dietterich}.

What is new here is the decomposition, not the accuracy: four arms rather
than with/without, a control that removes content while holding the
scaffolding byte-for-byte fixed, a ceiling compiled from the same content,
an artifact frozen and digest-pinned before any question was read, and
6{,}416 transcripts released.
